% Part: first-order-logic
% Chapter: completeness
% Section: complete-sets

% Definition of complete consistent sets. Properties of complete sets required
% for completeness proved are in provability.tex in the
% chapter on the proof system used.

\documentclass[../../../include/open-logic-section]{subfiles}

\begin{document}

\iftag{FOL}
      {\olfileid[id]{fol}{com}{ccs}}
      {\olfileid[id]{pl}{com}{ccs}}

\olsection{Himpunan \usetoken{P}{sentence} yang Lengkap dan Konsisten}

\begin{defn}[Himpunan lengkap]
\ollabel{def:complete-set} Suatu himpunan~$\Gamma$ dari !!{sentence}s
bersifat \emph{!!{complete}} jika dan hanya jika, untuk setiap
!!{sentence}~$!A$, berlaku $!A \in \Gamma$ atau $\lnot !A \in \Gamma$.
\end{defn}

\begin{explain}
Himpunan kalimat yang !!^{complete} tidak membiarkan pertanyaan apa pun
tanpa jawaban. Untuk setiap !!{sentence}~$!A$, $\Gamma$ ``menyatakan''
apakah $!A$ benar atau salah. Arti penting himpunan !!{complete} tidak
terbatas pada pembuktian teorema kelengkapan. Sebagai contoh, teori yang
!!{complete} dan dapat diaksiomakan selalu bersifat terputuskan.
\end{explain}

\begin{explain}
Himpunan konsisten yang !!^{complete} penting dalam pembuktian
kelengkapan karena kita dapat menjamin bahwa setiap himpunan konsisten
dari !!{sentence}s~$\Gamma$ termuat dalam suatu himpunan konsisten
$\Gamma^*$ yang !!a{complete}. Suatu himpunan konsisten yang
!!^a{complete} memuat, untuk setiap !!{sentence}~$!A$, $!A$ atau
negasinya $\lnot !A$, tetapi tidak keduanya. Hal ini berlaku khususnya bagi
\iftag{FOL}{!!{sentence}s atomik}{variabel proposisi}; maka, dari suatu
himpunan konsisten yang !!a{complete}\iftag{FOL}{ dalam bahasa yang
diperluas secara tepat dengan !!{constant}s}{}, kita dapat membangun
\iftag{FOL}{!!a{structure}}{!!a{valuation}} dengan
\iftag{FOL}{interpretasi !!{predicate}s}{nilai kebenaran yang ditetapkan
kepada variabel proposisi} didefinisikan menurut
\iftag{FOL}{!!{sentence}s atomik}{!!{propositional variable}s} mana yang
berada dalam~$\Gamma^*$. Selanjutnya, dapat dibuktikan bahwa
\iftag{FOL}{!!{structure}}{!!{valuation}} ini membuat semua
!!{sentence}s dalam~$\Gamma^*$ (dan karena itu juga semua kalimat
dalam~$\Gamma$) benar. Pembuktian fakta terakhir ini mengharuskan bahwa
$\lnot !A \in \Gamma^*$ jika dan hanya jika $!A \notin \Gamma^*$,
$(!A \lor !B) \in \Gamma^*$ jika dan hanya jika $!A \in \Gamma^*$
atau $!B \in \Gamma^*$, dan seterusnya.
\end{explain}

Dalam pembahasan berikut, kita akan sering menggunakan secara tersirat
sifat refleksivitas, monotonisitas, dan transitivitas $\Proves$ (lihat
\iftag{FOL}{%
  \tagrefs{prfSC/{fol:seq:ptn:sec},prfND/{fol:ntd:ptn:sec},prfAX/{fol:axd:ptn:sec},prfTab/{fol:tab:ptn:sec}}}{%
  \tagrefs{prfSC/{pl:seq:ptn:sec},prfND/{pl:ntd:ptn:sec},prfAX/{pl:axd:ptn:sec},prfTab/{pl:tab:ptn:sec}}}).

\begin{prop}
\ollabel{prop:ccs}
Andaikan $\Gamma$ !!{complete} dan konsisten. Maka:
\begin{enumerate}
\item \ollabel{prop:ccs-prov-in} Jika $\Gamma \Proves !A$, maka $!A \in
  \Gamma$.

\tagitem{prvAnd}{\ollabel{prop:ccs-and} $!A \land !B \in \Gamma$
  jika dan hanya jika $!A \in \Gamma$ dan $!B \in \Gamma$.}{}

\tagitem{prvOr}{\ollabel{prop:ccs-or} $!A \lor !B \in \Gamma$ jika dan
  hanya jika $!A \in \Gamma$ atau $!B \in \Gamma$.}{}

\tagitem{prvIf}{\ollabel{prop:ccs-if} $!A \lif !B \in \Gamma$ jika dan
  hanya jika $!A \notin \Gamma$ atau $!B \in \Gamma$.}{}
\end{enumerate}
\end{prop}

\begin{proof}
Untuk seluruh pembahasan berikut, andaikan bahwa $\Gamma$ !!{complete}
dan konsisten.
\begin{enumerate}
\item Jika $\Gamma \Proves !A$, maka $!A \in \Gamma$.

Andaikan $\Gamma \Proves !A$. Demi memperoleh kontradiksi, andaikan bahwa
$!A \notin \Gamma$. Karena $\Gamma$ !!{complete}, $\lnot !A \in \Gamma$.
Menurut
\iftag{FOL}{%
  \tagrefs{prfSC/{fol:seq:prv:prop:explicit-inc},
    prfND/{fol:ntd:prv:prop:explicit-inc},
    prfAX/{fol:axd:prv:prop:explicit-inc},
    prfTab/{fol:tab:prv:prop:explicit-inc}}}{%
  \tagrefs{prfSC/{pl:seq:prv:prop:explicit-inc},
    prfND/{pl:ntd:prv:prop:explicit-inc},
    prfAX/{pl:axd:prv:prop:explicit-inc},
    prfTab/{pl:tab:prv:prop:explicit-inc}}},
$\Gamma$ tak konsisten. Hal ini bertentangan dengan asumsi bahwa
$\Gamma$ konsisten. Jadi, tidak mungkin $!A \notin \Gamma$, sehingga
$!A \in \Gamma$.

\tagitem{defAnd}{}{%
\iftag{probAnd}{Latihan.}{%
$!A \land !B \in \Gamma$ jika dan hanya jika $!A \in \Gamma$ dan
$!B \in \Gamma$:

Untuk arah maju, andaikan $!A \land !B \in \Gamma$. Maka, menurut
\iftag{FOL}{%
  \tagrefs{prfSC/{fol:seq:ppr:prop:provability-land},
    prfND/{fol:ntd:ppr:prop:provability-land},
    prfAX/{fol:axd:ppr:prop:provability-land},
    prfTab/{fol:tab:ppr:prop:provability-land}}}{%
  \tagrefs{prfSC/{pl:seq:ppr:prop:provability-land},
    prfND/{pl:ntd:ppr:prop:provability-land},
    prfAX/{pl:axd:ppr:prop:provability-land},
    prfTab/{pl:tab:ppr:prop:provability-land}}}, butir~(1),
$\Gamma \Proves !A$ dan $\Gamma \Proves !B$. Menurut
\olref{prop:ccs-prov-in}, $!A \in \Gamma$ dan $!B \in \Gamma$, seperti
yang diperlukan.

Untuk arah sebaliknya, misalkan $!A \in \Gamma$ dan $!B \in
\Gamma$. Menurut
\iftag{FOL}{%
  \tagrefs{prfSC/{fol:seq:ppr:prop:provability-land},
    prfND/{fol:ntd:ppr:prop:provability-land},
    prfAX/{fol:axd:ppr:prop:provability-land},
    prfTab/{fol:tab:ppr:prop:provability-land}}}{%
  \tagrefs{prfSC/{pl:seq:ppr:prop:provability-land},
    prfND/{pl:ntd:ppr:prop:provability-land},
    prfAX/{pl:axd:ppr:prop:provability-land},
    prfTab/{pl:tab:ppr:prop:provability-land}}}, butir~(2),
$\Gamma \Proves !A \land !B$. Menurut \olref{prop:ccs-prov-in},
$!A \land !B \in \Gamma$.}}

\tagitem{defOr}{}{%
\iftag{probOr}{Latihan.}{%
Pertama, kita tunjukkan bahwa jika $!A \lor !B \in \Gamma$, maka
$!A \in \Gamma$ atau $!B \in \Gamma$. Andaikan $!A \lor !B \in
\Gamma$, tetapi $!A \notin \Gamma$ dan $!B \notin \Gamma$. Karena
$\Gamma$ !!{complete}, $\lnot !A \in \Gamma$ dan $\lnot !B \in
\Gamma$. Menurut
\iftag{FOL}{%
  \tagrefs{prfSC/{fol:seq:ppr:prop:provability-lor},
    prfND/{fol:ntd:ppr:prop:provability-lor},
    prfAX/{fol:axd:ppr:prop:provability-lor},
    prfTab/{fol:tab:ppr:prop:provability-lor}}}{%
  \tagrefs{prfSC/{pl:seq:ppr:prop:provability-lor},
    prfND/{pl:ntd:ppr:prop:provability-lor},
    prfAX/{pl:axd:ppr:prop:provability-lor},
    prfTab/{pl:tab:ppr:prop:provability-lor}}},
butir~(1), $\Gamma$ tak konsisten; ini suatu kontradiksi. Jadi,
$!A \in \Gamma$ atau $!B \in \Gamma$.

Untuk arah sebaliknya, andaikan $!A \in \Gamma$ atau $!B \in
\Gamma$. Menurut
\iftag{FOL}{%
  \tagrefs{prfSC/{fol:seq:ppr:prop:provability-lor},
    prfND/{fol:ntd:ppr:prop:provability-lor},
    prfAX/{fol:axd:ppr:prop:provability-lor},
    prfTab/{fol:tab:ppr:prop:provability-lor}}}{%
  \tagrefs{prfSC/{pl:seq:ppr:prop:provability-lor},
    prfND/{pl:ntd:ppr:prop:provability-lor},
    prfAX/{pl:axd:ppr:prop:provability-lor},
    prfTab/{pl:tab:ppr:prop:provability-lor}}}, butir~(2),
$\Gamma \Proves !A \lor !B$. Menurut \olref{prop:ccs-prov-in},
$!A \lor !B \in \Gamma$, seperti yang diperlukan.}}

\tagitem{defIf}{}{%
\iftag{probIf}{Latihan.}{%
Untuk arah maju, andaikan $!A \lif !B \in \Gamma$, lalu, demi memperoleh
kontradiksi, andaikan bahwa $!A \in \Gamma$ dan $!B \notin \Gamma$.
Berdasarkan asumsi-asumsi ini, $!A \lif !B \in \Gamma$ dan $!A \in \Gamma$.
Menurut
\iftag{FOL}{%
  \tagrefs{prfSC/{fol:seq:ppr:prop:provability-lif},
    prfND/{fol:ntd:ppr:prop:provability-lif},
    prfAX/{fol:axd:ppr:prop:provability-lif},
    prfTab/{fol:tab:ppr:prop:provability-lif}}}{%
  \tagrefs{prfSC/{pl:seq:ppr:prop:provability-lif},
    prfND/{pl:ntd:ppr:prop:provability-lif},
    prfAX/{pl:axd:ppr:prop:provability-lif},
    prfTab/{pl:tab:ppr:prop:provability-lif}}}, butir~(1),
$\Gamma \Proves !B$. Namun, menurut \olref{prop:ccs-prov-in}, hal ini
berarti $!B \in \Gamma$, bertentangan dengan asumsi bahwa $!B \notin
\Gamma$.

Untuk arah sebaliknya, pertama-tama tinjau kasus ketika $!A \notin
\Gamma$. Karena $\Gamma$ !!{complete}, $\lnot !A \in \Gamma$. Menurut
\iftag{FOL}{%
  \tagrefs{prfSC/{fol:seq:ppr:prop:provability-lif},
    prfND/{fol:ntd:ppr:prop:provability-lif},
    prfAX/{fol:axd:ppr:prop:provability-lif},
    prfTab/{fol:tab:ppr:prop:provability-lif}}}{%
  \tagrefs{prfSC/{pl:seq:ppr:prop:provability-lif},
    prfND/{pl:ntd:ppr:prop:provability-lif},
    prfAX/{pl:axd:ppr:prop:provability-lif},
    prfTab/{pl:tab:ppr:prop:provability-lif}}}, butir~(2),
$\Gamma \Proves !A \lif !B$. Sekali lagi, menurut
\olref{prop:ccs-prov-in}, kita memperoleh $!A \lif !B \in \Gamma$,
seperti yang diperlukan.

Sekarang tinjau kasus ketika $!B \in \Gamma$. Menurut
\iftag{FOL}{%
  \tagrefs{prfSC/{fol:seq:ppr:prop:provability-lif},
    prfND/{fol:ntd:ppr:prop:provability-lif},
    prfTab/{fol:tab:ppr:prop:provability-lif},
    prfAX/{fol:axd:ppr:prop:provability-lif}}}{%
  \tagrefs{prfSC/{pl:seq:ppr:prop:provability-lif},
    prfND/{pl:ntd:ppr:prop:provability-lif},
    prfTab/{pl:tab:ppr:prop:provability-lif},
    prfAX/{pl:axd:ppr:prop:provability-lif}}},
butir~(2) sekali lagi, $\Gamma \Proves !A \lif !B$. Menurut
\olref{prop:ccs-prov-in}, $!A \lif !B \in \Gamma$.}}
\end{enumerate}
\end{proof}

\tagprob{FOL}
\begin{prob}
Lengkapi pembuktian \olref[fol][com][ccs]{prop:ccs}.
\end{prob}
\tagendprob

\tagprob{notFOL}
\begin{prob}
Lengkapi pembuktian \olref[pl][com][ccs]{prop:ccs}.
\end{prob}
\tagendprob

\end{document}
